\documentclass[10pt,a4paper]{amsart}
\usepackage[margin=19mm]{geometry}
\usepackage{amsmath,amssymb,mathtools}
\usepackage[scr=rsfs,cal=euler]{mathalfa}
\usepackage{xfrac}
\usepackage[unicode,hidelinks,bookmarks=false]{hyperref}
\newcommand{\ie}{i.e.\ }
\newcommand{\cf}{cf.\ }
\newcommand{\resp}{respectively\ }
\newcommand{\Subsection}{\subsection}
\providecommand{\nobreakdash}{\nobreak\hbox{-}}
\setlength{\emergencystretch}{2em}
\multlinegap0pt
\usepackage{bm}
\usepackage{bbm}
\usepackage{dsfont}
\usepackage{xspace}
\usepackage{cancel}
\usepackage{mathtools}
\usepackage{amsthm}
\makeatletter
\@ifundefined{theorem}{\newtheorem{theorem}{Theorem}}{}
\makeatother
\makeatletter
\newcommand{\bbm}[1]{\left[\begin{matrix} #1 \end{matrix}\right]}
\expandafter\ifx\csname proof\endcsname\relax
\renewenvironment{proof}{\vspace{.1cm}\noindent{\sc Proof.}\hspace{0.10cm}\,\,}{$\hfill\Box$\vspace{.1cm}} 
\fi
\makeatother
\title[Neural network based control of unknown nonlinear systems vi]{Neural network based control of unknown nonlinear systems via contraction analysis}
\author{Hao Yin, Claudio De Persis, Bayu Jayawardhana, Santiago Sanchez Escalonilla Plaza}
\date{Source: arXiv:2505.16511v1 (CC BY 4.0)}
\begin{document}
\maketitle

\noindent\emph{Source and licence.} Neural network based control of unknown nonlinear systems via contraction analysis. Version arXiv:2505.16511v1 (CC BY 4.0), distributed under the Creative Commons Attribution 4.0 International licence, as recorded in the arXiv licence metadata for this exact version and shown on the abstract page https://arxiv.org/abs/2505.16511v1. Full rights notice: licenses/arxiv-2505.16511-CC-BY-4.0.txt.\par\smallskip

\noindent\textit{Editorial scope.} Counted unit: the Theorem whose source label is \texttt{theorem3} in the article (source0.tex lines 771--792, source label \texttt{theorem3}), delivered together with the article's own complete original proof of it (source0.tex lines 793--933, one \texttt{proof} environment of 141 lines). No mathematics is written, repaired or re-derived: statement and proof are the article's own text line for line. Required context retained uncounted: NN (lines 119--125); LEM1 (lines 183--199); NDL (lines 234--236); TH110 (lines 267--269); theorem1 (lines 374--390); theorem11 (lines 456--472); TH121P1 (lines 475--489); TH121P4 (lines 517--529); NNC10 (lines 615--617); NNC3 (lines 629--631); NNC11 (lines 641--643); CO31 (lines 674--676); theorem2 (lines 751--767); TH22 (lines 753--765). Every definition, display and label of those lines is retained unchanged. Omitted from this package and not redistributed: 1--118, 126--182, 200--233, 237--266, 270--373, 391--455, 473--474, 490--516, 530--614, 618--628, 632--640, 644--673, 677--750, 766--770, 934--1221. Nothing inside a retained excerpt is omitted or altered: every retained body is delivered exactly as the article's lines are written, so no omission receipt of any kind accompanies this package. Citations: the retained excerpts carry no citation command at all, so no bibliography entry of the article is transcribed and no reference is redistributed. Reference adaptation: none was needed and none was made; every reference inside the delivered unit resolves inside it, so no unresolved reference or citation is produced. Printed slips kept literal: no printed word, symbol, hypothesis or number of the counted statement or of its proof was altered. Layout and class: the article's own class is \texttt{ieeecolor} with packages including generic, amsmath, amssymb, amsfonts, algorithmic, graphicx, algorithm, hyperref, textcomp, mathrsfs, color, latexsym, bbm, dsfont; the wrapper is the lane's pinned \texttt{amsart} editorial wrapper at 10pt on A4, which restates the theorem-like environments the retained text needs and adds the article's own macro definitions that the retained text needs (\texttt{\textbackslash bbm}), with extra wrapper packages (bm, bbm, dsfont, xspace, cancel, mathtools). The delivered numbering is the unit's own. Assets: no retained span contains a figure environment, a graphics-inclusion command, a PDF-page inclusion, a table or a TikZ picture; no image file is copied into this package, the package declares no asset dependency and no image file is redistributed with it. Licence: version arXiv:2505.16511v1 of this article is distributed under the Creative Commons Attribution 4.0 International licence, as recorded in the arXiv licence metadata for this exact version and shown on the abstract page https://arxiv.org/abs/2505.16511v1. A full rights notice is delivered at licenses/arxiv-2505.16511-CC-BY-4.0.txt. Characters: every retained excerpt is pure ASCII, so the delivered engine, XeLaTeX with its default Latin Modern text font, reports no missing character.
\begin{equation}\label{NN}
 \begin{aligned}
&w^0=x,
\\&w^{i+1}=\phi^{i+1}(W^{i+1}w^{i}+b^{i+1}), i=0,...,k-1,
\\&Z_{nn}(x)=W^{k+1}w^{k}+b^{k+1},
\end{aligned}
\end{equation}

\begin{equation}\label{LEM1}
 \begin{aligned}
\begin{bmatrix}
w_1^{\phi}-w_2^{\phi}
 \\
\bar{w}_1^{\phi}-\bar{w}_2^{\phi}
\end{bmatrix}^\top \begin{bmatrix}
 -2Q_1^\top Q_2&Q_1^{\top}+Q_2^\top
 \\
\ast&-2I
\end{bmatrix}\begin{bmatrix}
w_1^{\phi}-w_2^{\phi}
 \\
\bar{w}_1^{\phi}-\bar{w}_2^{\phi}
\end{bmatrix}\geq 0, 
\end{aligned}
\end{equation}

 \begin{equation}\label{NDL}
\dot{x}_{nn}(t)=Ax_{nn}+Z_{nn}(x_{nn}), 
\end{equation}

\begin{equation}\label{TH110}
 \underline{p}I\leq P\leq \overline{p}I, 
\end{equation}

\begin{theorem}\label{theorem1}
Suppose there exist a positive definite matrix $P$, positive constants $\underline{p}$, $\overline{p}$, and $\gamma$ such that \eqref{TH110} and the following condition 
\begin{equation}\label{TH12}
\begin{aligned}
&U_1^\top
\begin{bmatrix}
     A^\top P+PA+\gamma P &PW^{k+1}\\
  \ast& 0
\end{bmatrix}U_1+\\&U_2^\top\begin{bmatrix}
 -2Q_1^\top Q_2&Q_1^\top+Q_2^\top
 \\
\ast&-2I
\end{bmatrix}U_2\leq 0, \quad\forall k\geq 2,
\end{aligned}
\end{equation}
hold, where $U_1=\bbm{I_{m_{k+1}\times m_{k+1}} &0_{m_{k+1}\times a}&0_{m_{k+1}\times m_{k}}\\0_{m_{k}\times m_{k+1}}&0_{m_{k}\times a}&I_{m_{k}\times m_{k}}}$, $U_2=\bbm{I_{m_{k+1}\times m_{k+1}} &0_{m_{k+1}\times a}&0_{m_{k+1}\times m_{k}}\\0_{a\times m_{k+1}}&I_{a\times a}&0_{a\times m_{k}}\\0_{a\times m_{k+1}}&I_{a\times a}&0_{a\times m_{k}}\\0_{m_{k}\times m_{k+1}}&0_{m_{k}\times a}&I_{m_{k}\times m_{k}}}$, $a=\sum_{i=1}^{k-1}m_i$, $W^{k+1}$, $Q_1$, $Q_2$ are given in \eqref{NN}, \eqref{LEM1}, respectively. Then the NODE system \eqref{NDL} is contractive in $\mathcal{C}$.
\end{theorem}

\begin{theorem}\label{theorem11}
Suppose there exist a positive definite matrix $P$, positive constants $\underline{p}$, $\overline{p}$, and $\mu$, such that \eqref{TH110} and the following condition
\begin{equation}\label{TH121}
\begin{aligned}
&\bar{U}_1^\top
\bbm{-\mu I&P&0\\
   \ast&  A^\top P+PA &PW^{k+1}\\
  \ast& \ast & 0
}\bar{U}_1+\\&\bar{U}_2^\top\bbm{0&0&0\\\ast&
 -2Q_1^\top Q_2&Q_1^\top+Q_2^\top
 \\\ast&
\ast&-2I
}\bar{U}_2\leq 0, \quad\forall k\geq 2,
\end{aligned}
\end{equation}
hold, where \\$\bar{U}_1=\bbm{I_{m_{k+1}\times m_{k+1}} &0_{m_{k+1}\times m_{k+1}}&0_{m_{k+1}\times a} &0_{m_{k+1}\times m_{k}}\\0_{m_{k+1}\times m_{k+1}}&I_{m_{k+1}\times m_{k+1}}&0_{m_{k+1}\times a}&0_{m_{k+1}\times m_{k}}\\0_{m_{k}\times m_{k+1}}&0_{m_{k}\times m_{k+1}}&0_{m_{k}\times a}&I_{m_{k}\times m_{k}}}$, $\bar{U}_2=\bbm{I_{m_{k+1}\times m_{k+1}} &0_{m_{k+1}\times m_{k+1}}&0_{m_{k+1}\times a}&0_{m_{k+1}\times m_{k}}\\0_{m_{k+1}\times m_{k+1}}&I_{m_{k+1}\times m_{k+1}} &0_{m_{k+1}\times a}&0_{m_{k+1}\times m_{k}}\\0_{a\times m_{k+1}}&0_{a\times m_{k+1}}&I_{a\times a}&0_{a\times m_{k}}\\0_{a\times m_{k+1}}&0_{a\times m_{k+1}}&I_{a\times a}&0_{a\times m_{k}}\\0_{m_{k}\times m_{k+1}}&0_{m_{k}\times m_{k+1}}&0_{m_{k}\times a}&I_{m_{k}\times m_{k}}}$, $a=\sum_{i=1}^{k-1}m_i$, $W^{k+1}$, $Q_1$, $Q_2$ are given in \eqref{NN}, \eqref{LEM1}, respectively. Then the NODE system \eqref{NDL} is contractive in $\mathcal{C}$.
\end{theorem}

 \begin{equation}\label{TH121P1}
\begin{aligned}
&\bar{U}_1^\top
\bbm{-\mu I&P&0\\
   \ast&  A^\top P+PA &PW^{k+1}\\
  \ast& \ast & 0
}\bar{U}_1\\&=\bbm{I&0\\0&U^\top_1}\bbm{-\mu I&P&0\\
   \ast&  A^\top P+PA &PW^{k+1}\\
  \ast& \ast & 0
}\bbm{I&0\\0&U_1}\\&=
\bbm{-\mu I&\bbm{P&0}\\U_1^\top \bbm{P\\0}&U_1^\top\bbm{A^\top P+PA &PW^{k+1}\\
 \ast& 0}}\bbm{I_{n\times n}&0\\0&U_1}\\&=\bbm{-\mu I&\bbm{P&0}U_1\\\ast&U_1^\top\bbm{A^\top P+PA &PW^{k+1}\\
  \ast & 0}U_1}.
\end{aligned}
\end{equation}

\begin{equation}\label{TH121P4}
\begin{aligned}
&U_1^\top\bbm{\frac{1}{\mu}P^2&0\\0&0}U_1+\bar{L}=\\&U_1^\top
\begin{bmatrix}
     A^\top P+PA+\frac{1}{\mu}P^2 &PW^{k+1}\\
   \ast & 0
\end{bmatrix}U_1+\\&U_2^\top\begin{bmatrix}
 -2Q_1^\top Q_2&Q_1^\top+Q_2^\top
 \\
\ast&-2I
\end{bmatrix}U_2\leq 0.
\end{aligned}
\end{equation}

\begin{equation}\label{NNC10}
u(x)=Hx+u_{nn}(x),
\end{equation}

\begin{equation}\label{NNC3}
  \dot{x}=(A+gH)x+(W^{k+1}+gW_u^{k+1})w^k+b^{k+1}+gb_u^{k+1}.
\end{equation}

\begin{equation}\label{NNC11}
u(x)=Hx+W_{u\min}^{k+1}w^k+b_u^{k+1}.
\end{equation}

\begin{equation}\label{CO31}
 \underline{s}I\leq S\leq \overline{s}I, 
\end{equation}

\begin{theorem}\label{theorem2}
Suppose there exist a positive definite matrix $P$, positive constants $\underline{p}$, $\overline{p}$, $\gamma$, and a matrix $H$, such that \eqref{TH110} and the following condition 
\begin{equation}\label{TH22}
\begin{aligned}
&U_1^\top
\begin{bmatrix}
     (A+gH)^\top P+P(A+gH)+\gamma P &PW_{\min}\\
   \ast & 0
\end{bmatrix}U_1+\\&U_2^\top\begin{bmatrix}
 -2Q_1^\top Q_2&Q_1^\top+Q_2^\top
 \\
\ast&-2I
\end{bmatrix}U_2\leq 0,
\end{aligned}
\end{equation}
hold, where $U_1$, $U_2$, $W^{k+1}$, $Q_1$, $Q_2$ are given in Theorem \ref{theorem1}. Then, the learned affine system \eqref{NNC3} is contractive in $\mathcal{C}$ under the controller specified by \eqref{NNC11}.
\end{theorem}

\begin{theorem}\label{theorem3}
Suppose there exist a fixed positive constant $\underline{s}$, positive definite matrices $S$, positive constants $\mu$,  $\overline{s}$, and a matrix $Y$, such that \eqref{CO31} and the following condition
\begin{equation}\label{TH32}
\begin{aligned}
&\bar{U}_1^\top
\bbm{
-\mu I & S & 0 \\
\ast &  (AS+gY)^\top + AS + gY & W_{\min} \\
\ast & \ast & 0
}
\bar{U}_1 + \\
&\bar{U}_2^\top
\bbm{
0 & 0 & 0 \\
\ast & \bar{X} &\tilde{X} \\
\ast & \ast & -2I
}
\bar{U}_2 \leq 0, \quad \forall k \geq 2, 
\end{aligned}
\end{equation}
hold, where $\bar{X}=\bbm{-2\underline{s} \underline{\lambda}S&0\\0&-2Q^{a\top}_1 Q^a_2}$, $\underline{\lambda}$ is the minimum eigenvalue of $Q^{1\top}_1 Q^1_2$,  $\tilde{X}=\bbm{S(Q^1_1+Q^1_2)^\top&0\\0&(Q^a_1+Q^a_2)^\top}$, $Q^a_1=\mathrm{BD}(K_1^2W^2, \dots, K_1^{k}W^{k})$, $Q^a_2=\mathrm{BD}(K_2^2W^2, \dots, K_2^{k}W^{k})$, $\bar{U}_1$, $\bar{U}_2$, $W^{k+1}$ are given in Theorem \ref{theorem11}. Then, the learned affine system \eqref{NNC3} is contractive in $\mathcal{C}$ under the controller specified by \eqref{NNC10} with $H=YS^{-1}$.
\end{theorem}

\begin{proof}
Observe that $\bar{U}_1=\bbm{I_{m_{k+1}\times m_{k+1}}&0\\0&U_1}$, $\bar{U}_2=\bbm{I_{m_{k+1}\times m_{k+1}}&0\\0&U_2}$. Following a similar approach as \eqref{TH121P1}-\eqref{TH121P4} in the proof of Theorem \ref{theorem11}, but substituting $A$ and $P$ with $AS + gY$ and $S$, respectively, we can rewrite \eqref{TH32} as
\begin{equation}\label{TH3P1}
\begin{aligned}
&U_1^\top
\bbm{
(AS+gY)^\top + AS + gY+\frac{1}{\mu}S^2 & W_{\min} \\
\ast & 0
}
U_1 + \\
&U_2^\top
\bbm{
\bar{X} &\tilde{X} \\
\ast & -2I
}
U_2 \leq 0, 
\end{aligned}
\end{equation}
We can express the left hand side of \eqref{TH3P1} in the following form:
\begin{equation}\label{TH3P40}
\begin{aligned}
&\Big(\bbm{
S&0\\0&I
}U_1\Big)^\top
\bbm{
X & S^{-1}W_{\min} \\
\ast & 0
}
\bbm{
S&0\\0&I
}U_1 + \\
&\Big(\bbm{
S&0&0&0\\\ast&I&0&0\\\ast&\ast&I&0\\\ast&\ast&\ast&I
}U_2\Big)^\top\\&\times
\bbm{\bbm{-2\underline{s}\underline{\lambda}S^{-1}&0\\0&-2Q^{a\top}_1 Q^a_2}
 &*\\
\bbm{(Q^1_1+Q^1_2)^\top&0\\0&(Q^a_1+Q^a_2)^\top}^\top  & -2I
}
\\&\times\bbm{
S&0&0&0\\\ast&I&0&0\\\ast&\ast&I&0\\\ast&\ast&\ast&I
}U_2, 
\end{aligned}
\end{equation}
where $X=(S^{-1}A + S^{-1}gYS^{-1})^\top + S^{-1}A + S^{-1}gYS^{-1}+\frac{1}{\mu}I$.

Note that
\begin{equation}\label{TH3P2}
\begin{aligned}
&U_1\bbm{
S&0&0\\\ast&I&0\\\ast&\ast&I
}=\bbm{
S&0\\0&I
}U_1,
\end{aligned}
\end{equation}
and
\begin{equation}\label{TH3P3}
\begin{aligned}
&U_2\bbm{
S&0&0\\\ast&I&0\\\ast&\ast&I
}=\bbm{
S&0&0&0\\\ast&I&0&0\\\ast&\ast&I&0\\\ast&\ast&\ast&I
}U_2,
\end{aligned}
\end{equation}
Using \eqref{TH3P2}, \eqref{TH3P3}, we can rewrite \eqref{TH3P4} as:
\begin{equation}\label{TH3P4}
\begin{aligned}
&\bbm{
S&0&0\\\ast&I&0\\\ast&\ast&I
}\Big(U_1^\top
\bbm{
X & S^{-1}W_{\min} \\
W_{\min}^{\top}S^{-1} & 0
}
U_1 +\\
& U_2^\top
\bbm{\bbm{-2\underline{s}\underline{\lambda}S^{-1}&0\\0&-2Q^{a\top}_1 Q^a_2}
 &\ast\\
\bbm{(Q^1_1+Q^1_2)^\top&0\\0&(Q^a_1+Q^a_2)^\top}^\top  & -2I
}
U_2\Big)\\&\times\bbm{
S&0&0\\\ast&I&0\\\ast&\ast&I
}.
\end{aligned}
\end{equation}
Using $P=S^{-1}$, $H=YS^{-1}$, we can rewrite \eqref{TH3P4} as
\begin{equation}\label{TH3P5}
\begin{aligned}
&\bbm{
P^{-1}&0&0\\\ast&I&0\\\ast&\ast&I
}\times\\&\Big(U_1^\top
\bbm{
(A+gH)^\top P+P(A+gH)+\frac{1}{\mu}I & PW_{\min} \\
\ast & 0
}
U_1 \\
&+ U_2^\top
\bbm{\bbm{-2\underline{s}\underline{\lambda}P&0\\0&-2Q^{a\top}_1 Q^a_2}
 &\ast\\
\bbm{(Q^1_1+Q^1_2)^\top&0\\0&(Q^a_1+Q^a_2)^\top}^\top  & -2I
}
U_2\Big)\\&\times\bbm{
P^{-1}&0&0\\\ast&I&0\\\ast&\ast&I
}.
\end{aligned}
\end{equation}
Since $\bbm{
P^{-1}&0&0\\\ast&I&0\\\ast&\ast&I
}$ is invertible, using \eqref{TH3P5}, the inequality \eqref{TH3P1} is equivalent to
\begin{equation}\label{TH3P6}
\begin{aligned}
&U_1^\top
\bbm{
(A+gH)^\top P+P(A+gH)+\frac{1}{\mu}I & PW_{\min} \\
\ast & 0
}
U_1 +\\
& U_2^\top
\bbm{\bbm{-2\underline{s}\underline{\lambda}P&0\\0&-2Q^{a\top}_1 Q^a_2}
 &\ast\\
\bbm{(Q^1_1+Q^1_2)^\top&0\\0&(Q^a_1+Q^a_2)^\top}^\top  & -2I
}
U_2\leq 0.
\end{aligned}
\end{equation}
Using $\underline{s}P\geq I$, $\gamma=\frac{\underline{s}}{\mu}$, we have 
\begin{equation}\label{TH3P7}
\begin{aligned}
&U_1^\top\bbm{\gamma P&0\\0&0}U_1=\bbm{\gamma P&0&0\\\ast&0&0\\\ast&\ast&0}\leq \bbm{\frac{1}{\mu}I&0&0\\\ast&0&0\\\ast&\ast&0}\\&=U_1^\top\bbm{\frac{1}{\mu}I&0\\0&0}U_1,
\end{aligned}
\end{equation}
and
\begin{equation}\label{TH3P8}
\begin{aligned}
&U_2^\top\bbm{-2Q_1^{1\top}Q_2^1&0&0&0\\\ast&0&0&0\\\ast&\ast&0&0\\\ast&\ast&\ast&0}U_2=\bbm{-2Q_1^{1\top}Q_2^1&0&0&0\\\ast&0&0&0\\\ast&\ast&0&0\\\ast&\ast&\ast&0}\\&\leq \bbm{-2\underline{s}\underline{\lambda}P&0&0&0\\\ast&0&0&0\\\ast&\ast&0&0\\\ast&\ast&\ast&0}=U_2^\top\bbm{-2\underline{s}\underline{\lambda}P&0&0&0\\\ast&0&0&0\\\ast&\ast&0&0\\\ast&\ast&\ast&0}U_2.
\end{aligned}
\end{equation}
Then inequity \eqref{TH22} in Theorem \ref{theorem1} is satisfied by applying inequities \eqref{TH3P6}, \eqref{TH3P7}, and \eqref{TH3P8}. The conditions in Theorem \ref{theorem2} are met and the
conclusion of the theorem follows directly. 
\end{proof}

\end{document}
